\begin{claimboundary}
長時間・tool-usingな実行能力が高まるほど、安全性評価はモデルの付録ではなく実行系の設計条件になる。GPT-5.3-CodexのSystem CardはPreparedness Frameworkのcybersecurity評価を適用している。ただし、その分類はSystem Cardが記載したthresholdとprecautionの文脈で読む必要があり、『危険な能力が確立した』といった強い表現へ置き換えることはできない。2月の技術的ポイントは、agentic capabilityの拡張と同時に、その実行可能性をどの境界で制御するかが同じ製品系列の論点として現れたことにある。\autocite{src-15823ae5556e,src-4a841ea7ec7c,src-9aba7995b6a5,src-dd64fe3084a4}
\end{claimboundary}

\begin{claimboundary}
性能比較には明確な境界を置く。GPT-5.3-Codexの速度向上やSparkのthroughput、Claude 4.6系のbenchmark・比較性能はいずれも各社の発表に基づくvendor claimであり、本稿で独立再現した値ではない。モデルのrelease、対象task、preview/betaといった一次資料で固定できる事実と、優位性を示す比較値は同じ確度で扱わない。agentic engineeringの方向性はbenchmarkの順位ではなく、製品とモデルが実際にどの実行surfaceへ拡張されたかから読む。\autocite{src-15823ae5556e,src-4a841ea7ec7c,src-9aba7995b6a5,src-dd64fe3084a4,src-f6aa679babd7}
\end{claimboundary}

2月のagentic engineeringには二つの時間軸が同居している。一つは、複数agentやlong-running research、computer useのようにAIへ長い仕事を任せる方向。もう一つは、Sparkのように人間とのrealtimeなcoding loopを短くする方向である。自律時間を延ばすことと対話遅延を縮めることは一見逆向きだが、どちらも『AIを実際の作業工程へ埋め込む』ための最適化である。3月以降を見る際にも、単純なmodel scoreより、この二つの時間軸と安全な制御面がどこまで実装へ落ちるかが重要な観察点になる。\autocite{src-15823ae5556e,src-4a841ea7ec7c,src-9aba7995b6a5,src-dd64fe3084a4,src-f6aa679babd7}
